\section{Constructive continuum sector, regulator removal, and physical identifiability}
\label{sec:constructive-continuum}

The finite-regulator theorem proves exact quantum identities before a continuum limit is taken.  This section supplies a nonperturbative continuum construction for the positive quadratic physical sector of the action-lifted Recognition theory.  It proves regulator removal, reflection positivity, anomaly freedom of the bosonic measure, and unitary Lorentzian reconstruction under explicit spectral hypotheses.  It then gives a target-faithful criterion for identifying the Recognition metric and connection from propagation, scale, and transport data.

\subsection{Positive physical Hessian datum}

Let $K$ be a real separable Hilbert space representing gauge-fixed, constraint-projected, target-faithful physical fluctuations about a stationary Euclidean background.  Let
\[
\Omega:\mathcal D(\Omega)\subset K\longrightarrow K
\]
be positive and self-adjoint with compact resolvent and a strict gap
\begin{equation}
\Omega\ge m\id,
\qquad m>0.
\label{eq:omega-gap}
\end{equation}
Define the smooth spectral test space
\[
\mathscr S_\Omega:=\bigcap_{r\ge0}\mathcal D(\Omega^r)
\]
with its graph seminorms, and assume that this Fr\'{e}chet space is nuclear.  Set
\[
\mathscr E
:=\mathscr S(\R_\tau)\widehat\otimes\mathscr S_\Omega,
\qquad
\mathscr E'
:=\text{its continuous dual}.
\]
The Euclidean quadratic operator is
\[
L_E=-\partial_\tau^2+\Omega^2.
\]
For $f,g\in\mathscr E$, define the covariance form
\begin{equation}
Q(f,g)
:=\int_{\R}\int_{\R}
\left\langle
f(\tau),
\frac{e^{-\abs{\tau-s}\Omega}}{2\Omega}
 g(s)
\right\rangle_K
\,\dd\tau\,\dd s.
\label{eq:continuum-covariance}
\end{equation}
Functional calculus and \eqref{eq:omega-gap} make $Q$ continuous, symmetric, and nonnegative.

Let
\[
P_N:=\mathbf 1_{[m,N]}(\Omega).
\]
Compact resolvent implies that $P_NK$ is finite dimensional.  The spatial spectral-cut covariance is
\begin{equation}
Q_N(f,g):=Q(P_Nf,P_Ng).
\label{eq:cutoff-covariance}
\end{equation}

\begin{theorem}[Continuum Gaussian Recognition measure and regulator removal]
\label{thm:continuum-gaussian}
There exist unique centered Gaussian probability measures $\mu_N$ and $\mu$ on $\mathscr E'$ with characteristic functionals
\begin{align}
\widehat\mu_N(f)
&=\exp\!\left[-\frac12Q_N(f,f)\right],
\label{eq:cutoff-characteristic}\\
\widehat\mu(f)
&=\exp\!\left[-\frac12Q(f,f)\right].
\label{eq:continuum-characteristic}
\end{align}
As $N\to\infty$,
\begin{equation}
\mu_N\Longrightarrow\mu
\label{eq:measure-convergence}
\end{equation}
weakly on $\mathscr E'$.  Every cylindrical moment converges, and the limiting Schwinger functions are the Wick functions generated by $Q$.  The same limit is obtained from any increasing sequence of finite-rank spectral projections commuting with $\Omega$ and converging strongly to the identity.  No counterterm is required in this positive quadratic sector.
\end{theorem}

\begin{proof}
The forms $Q_N$ and $Q$ are continuous and nonnegative on the nuclear space $\mathscr E$.  Their exponentials in \eqref{eq:cutoff-characteristic}--\eqref{eq:continuum-characteristic} are continuous positive-definite functionals equal to one at the origin.  The Bochner--Minlos theorem therefore gives the unique Gaussian measures \cite{Minlos1959}.

Because $P_N\to\id$ strongly, $P_N$ commutes with every bounded Borel function of $\Omega$, and $(2\Omega)^{-1}$ is bounded by $(2m)^{-1}$, dominated convergence in \eqref{eq:continuum-covariance} gives
\[
Q_N(f,f)\longrightarrow Q(f,f)
\quad\text{for every }f\in\mathscr E.
\]
Thus $\widehat\mu_N(f)\to\widehat\mu(f)$ pointwise.  The continuity at the origin is uniform on a sufficiently small test-space neighbourhood because $0\le Q_N(f,f)\le Q(f,f)$.  The continuity theorem for probability measures on nuclear duals gives \eqref{eq:measure-convergence}.

Gaussian cylindrical moments are finite sums of products of two-point covariances.  Their convergence follows from $Q_N(f,g)\to Q(f,g)$.  The proof uses only strong convergence, commutation with $\Omega$, and finite rank, so any such nested regulator has the same limit.  Since the unmodified covariances converge, no mass, wave-function, or coupling counterterm is introduced in this quadratic sector.
\end{proof}

\begin{corollary}[Gaussian continuum completion]
\label{cor:gaussian-uv}
The positive quadratic physical Recognition sector has a regulator-independent continuum measure and all of its Schwinger functions define continuous generalized fields on the nuclear dual $\mathscr E'$.  This is a nonperturbative continuum and ultraviolet completion of the linearized sector in the precise sense that every spectral regulator is removed without counterterms and all cylindrical correlations converge.  It does not assert ultraviolet completion of the nonlinear metric--gauge interaction.
\end{corollary}

\subsection{Reflection positivity}

Let Euclidean time reflection act by
\[
(\theta f)(\tau):=f(-\tau),
\]
and let $\mathscr E_+$ consist of test functions supported in $\tau>0$.  Define
\begin{equation}
Vf
:=\int_0^\infty
(2\Omega)^{-1/2}e^{-\tau\Omega}f(\tau)\,\dd\tau
\in K.
\label{eq:os-one-particle-map}
\end{equation}

\begin{lemma}[Reflection factorization]
\label{lem:reflection-factorization}
For $f,g\in\mathscr E_+$,
\begin{equation}
Q(\theta f,g)=\langle Vf,Vg\rangle_K.
\label{eq:reflection-factorization}
\end{equation}
The same identity holds with $Q_N$ and $V_N:=P_NV$.
\end{lemma}

\begin{proof}
The support assumptions place the reflected variable at negative time and the unreflected variable at positive time.  Hence the absolute value in \eqref{eq:continuum-covariance} becomes the sum of their positive time parameters.  Functional calculus gives
\[
\frac{e^{-(\tau+s)\Omega}}{2\Omega}
=(2\Omega)^{-1/2}e^{-\tau\Omega}
 e^{-s\Omega}(2\Omega)^{-1/2}.
\]
Fubini's theorem yields \eqref{eq:reflection-factorization}.  Inserting $P_N$ gives the cutoff identity.
\end{proof}

For a field $\phi\in\mathscr E'$, write $\phi(f)$ for evaluation.  Let $\mathcal C_+$ be the linear span of the positive-time exponential functionals
\[
F(\phi)=\sum_{j=1}^r c_j e^{i\phi(f_j)},
\qquad f_j\in\mathscr E_+.
\]

\begin{theorem}[Osterwalder--Schrader positivity]
\label{thm:os-positivity}
For every $F\in\mathcal C_+$,
\begin{equation}
\int_{\mathscr E'}
\overline{F(\theta\phi)}F(\phi)\,\dd\mu(\phi)
\ge0.
\label{eq:os-positivity}
\end{equation}
Every cutoff measure $\mu_N$ satisfies the same inequality, and reflection positivity is preserved by the regulator limit.
\end{theorem}

\begin{proof}
For $F=\sum_jc_je^{i\phi(f_j)}$, Gaussian evaluation and reflection invariance give
\begin{align*}
&\int\overline{F(\theta\phi)}F(\phi)\,\dd\mu(\phi)\\
&\quad=\sum_{j,k}\overline{c_j}c_k
\exp\!\left[-\frac12Q(f_j,f_j)-\frac12Q(f_k,f_k)
+Q(\theta f_j,f_k)\right].
\end{align*}
Set $d_j=c_j\exp[-Q(f_j,f_j)/2]$.  By Lemma~\ref{lem:reflection-factorization}, the remaining matrix is
\[
\exp\langle Vf_j,Vf_k\rangle_K.
\]
It is positive semidefinite because
\[
\exp\langle u,v\rangle
=\sum_{n=0}^\infty\frac1{n!}
\langle u^{\otimes n},v^{\otimes n}\rangle,
\]
so it is the Gram kernel of bosonic exponential vectors.  This proves \eqref{eq:os-positivity}.  The cutoff proof is identical with $V_N$, and the continuum statement also follows from convergence of the finite matrices.  This is the free-field reflection-positive construction underlying the Osterwalder--Schrader reconstruction theorem \cite{OS1973,OS1975}.
\end{proof}

\subsection{Lorentzian unitary reconstruction}

Let $\mathcal N$ be the null space of the sesquilinear form in \eqref{eq:os-positivity}.  Define
\[
\mathcal H_{\mathrm{OS}}
:=\overline{\mathcal C_+/\mathcal N}.
\]
For $t\ge0$, translate a positive-time test function away from the reflection plane by
\[
(\tau_tf)(s):=f(s-t),
\]
whenever the translated support remains positive.

\begin{theorem}[Positive Hamiltonian and Lorentzian unitarity]
\label{thm:os-unitarity}
The Osterwalder--Schrader Hilbert space is canonically isomorphic to the symmetric Fock space $\Gamma_s(K_{\mathrm{phys}}^{\mathbb C})$ over the complexification of
\[
K_{\mathrm{phys}}:=\overline{\Ran V}\subseteq K.
\]
More precisely, the class of $e^{i\phi(f)}$ is mapped to
\[
\exp\!\left[-\frac12Q(f,f)\right]\operatorname{Exp}(Vf),
\]
where $\operatorname{Exp}$ is the bosonic exponential vector.  Under this identification, positive Euclidean time translation is
\begin{equation}
T_t=\Gamma_s(e^{-t\Omega})=e^{-tH},
\qquad
H=\dd\Gamma(\Omega)\ge0.
\label{eq:os-semigroup}
\end{equation}
Consequently,
\begin{equation}
U_t=e^{-itH}
\label{eq:lorentzian-unitary}
\end{equation}
is a strongly continuous unitary Lorentzian time evolution.  The same reconstruction is obtained as the limit of the finite-mode theories.
\end{theorem}

\begin{proof}
The proof of Theorem~\ref{thm:os-positivity} shows that the displayed weighted exponential vectors have exactly the Osterwalder--Schrader Gram matrix.  Their span is dense in the symmetric Fock space over $K_{\mathrm{phys}}^{\mathbb C}$, so the map extends by completion to a unitary identification.  From \eqref{eq:os-one-particle-map},
\[
V(\tau_tf)=e^{-t\Omega}Vf.
\]
Second quantization therefore gives \eqref{eq:os-semigroup}.  Since $\Omega\ge0$, $H=\dd\Gamma(\Omega)$ is self-adjoint and nonnegative.  Stone's theorem then gives the unitary group \eqref{eq:lorentzian-unitary}.  Strong convergence of $P_N$ gives strong convergence of the one-particle semigroups and hence convergence on the dense set of finite-particle vectors.
\end{proof}

\subsection{Bosonic anomaly cancellation}

Let a compact group $G$ act orthogonally on $K$ through $U:G\to\mathcal O(K)$.  Assume
\begin{equation}
U_g\Omega=\Omega U_g
\quad\text{for every }g\in G.
\label{eq:symmetry-commutes-omega}
\end{equation}
The action extends pointwise to $\mathscr E$ and dually to $\mathscr E'$.

\begin{theorem}[Exact bosonic anomaly freedom]
\label{thm:bosonic-anomaly-free}
Under \eqref{eq:symmetry-commutes-omega},
\begin{equation}
(U_g)_*\mu_N=\mu_N,
\qquad
(U_g)_*\mu=\mu
\label{eq:measure-invariance}
\end{equation}
for every $g\in G$.  Hence the bosonic spectral regulators and their continuum limit have no measure Jacobian anomaly.  The reconstructed symmetry is unitary on $\mathcal H_{\mathrm{OS}}$ and commutes with $H$.  For every differentiable cylindrical functional $F$ and infinitesimal generator $X$,
\begin{equation}
\int \mathcal L_XF\,\dd\mu=0.
\label{eq:bosonic-ward}
\end{equation}
\end{theorem}

\begin{proof}
Commutation with $\Omega$ gives $Q(U_gf,U_gg)=Q(f,g)$ and also makes $U_g$ commute with $P_N$.  Therefore the characteristic functionals of the pushforward measures equal those in \eqref{eq:cutoff-characteristic}--\eqref{eq:continuum-characteristic}; uniqueness of Gaussian measures proves \eqref{eq:measure-invariance}.  Differentiating the exact invariance at the identity gives \eqref{eq:bosonic-ward}.  The one-particle action preserves $K_{\mathrm{phys}}$, commutes with $e^{-t\Omega}$, and second quantizes to a unitary symmetry commuting with $H$.
\end{proof}

\begin{remark}[Scope of anomaly cancellation]
The theorem closes the anomaly question for the bosonic positive quadratic sector and for symmetry-preserving spectral regulators.  It does not decide chiral-fermion anomalies, nonlinear ghost determinants, or gravitational anomalies of a larger field content.
\end{remark}

\subsection{Application to linearized Recognition gravity}

\begin{corollary}[Constructive physical-mode Recognition sector]
\label{cor:constructive-physical-sector}
Let $(g_0,A_0,\Phi_0)$ be a stationary Euclidean background of the action-lifted Recognition equations.  Suppose that gauge fixing, constraint reduction, and target-faithful projection produce a real physical Hilbert space $K$ on which the quadratic Hessian has the ultrastatic form
\[
L_E=-\partial_\tau^2+\Omega^2
\]
with the hypotheses above.  Then the linearized physical metric, gauge, and Recognition-field modes have:
\begin{enumerate}[label=(\roman*)]
\item a regulator-independent continuum Gaussian measure;
\item spectral-cutoff removal with convergence of all cylindrical Schwinger functions;
\item Osterwalder--Schrader reflection positivity;
\item a positive reconstructed Hamiltonian and unitary Lorentzian evolution;
\item exact bosonic anomaly freedom for every compact symmetry commuting with the physical Hessian;
\item a nonperturbative Gaussian ultraviolet completion.
\end{enumerate}
\end{corollary}

\begin{remark}[Why positivity is load-bearing]
The corollary applies after projection to a positive physical Hessian.  It does not conceal a negative conformal mode, an unremoved gauge direction, or a ghost instability inside the word ``regulator.''  If the reduced Hessian is not positive, reflection positivity and the probability-measure construction must be re-established by a different argument.
\end{remark}

\section{Target-faithful identification of the Recognition spacetime geometry}
\label{sec:metric-identification}

A mathematical theory cannot manufacture experimental data.  It can, however, prove that a declared set of observations is sufficient to identify the geometric target up to its lawful symmetries.  The Recognition Kernel language makes this distinction exact.

\subsection{Metric identification from propagation and scale}

For a Lorentzian metric $g$, let $\mathcal N_g(x)$ denote its null cone in $T_xM$.  Let $\mathrm{vol}_g$ be its positive volume density.

\begin{theorem}[Null-cone and scale identification]
\label{thm:null-cone-identification}
Let $M$ be connected of dimension $n\ge3$, and let $g$ and $\widetilde g$ be time-oriented Lorentzian metrics.  If
\begin{equation}
\mathcal N_g(x)=\mathcal N_{\widetilde g}(x)
\quad\text{for every }x\in M,
\label{eq:same-null-cones}
\end{equation}
then there is a smooth positive function $\Omega_c$ such that
\begin{equation}
\widetilde g=\Omega_c^2g.
\label{eq:conformal-identification}
\end{equation}
The conformal factor is fixed, and hence $\widetilde g=g$, under either of the following pointwise scale calibrations:
\begin{enumerate}[label=(\alph*)]
\item $\mathrm{vol}_{\widetilde g}=\mathrm{vol}_g$;
\item there is a smooth nowhere-null vector field $T$ such that $\widetilde g(T,T)=g(T,T)$.
\end{enumerate}
More generally, if the observations are related by a diffeomorphism $\psi$, the conclusion is $g=\psi^*\widetilde g$.
\end{theorem}

\begin{proof}
At each point, two nondegenerate Lorentzian quadratic forms in dimension at least three with the same null cone are positive scalar multiples of one another.  To see this, choose a $g$-orthonormal basis with quadratic form $-t^2+\abs{x}^2$.  Evaluating $\widetilde g$ on all vectors $(1,u)$ with $\abs u=1$ forces the mixed coefficients to vanish and the spatial block to be a common scalar multiple of the identity; the time coefficient is then the negative of the same scalar.  Smoothness and time orientation make the scalar a smooth positive function, giving \eqref{eq:conformal-identification}.  Volume densities transform as
\[
\mathrm{vol}_{\widetilde g}=\Omega_c^n\mathrm{vol}_g,
\]
so calibration (a) and positivity imply $\Omega_c=1$.  Under calibration (b),
\[
\widetilde g(T,T)=\Omega_c^2g(T,T)=g(T,T),
\]
and the nowhere-null hypothesis again gives $\Omega_c=1$.  Pulling back by $\psi$ reduces the diffeomorphic statement to the preceding case.  The causal determination of conformal structure is the geometric principle formalized in \cite{Malament1977}.
\end{proof}

\begin{corollary}[Principal-symbol identification]
If two normally hyperbolic Recognition operators have the same characteristic covectors and either the same calibrated volume density or one common pointwise proper-time/length scale along a nowhere-null field, then their Lorentzian metrics agree up to diffeomorphism.  Thus high-frequency propagation determines the conformal metric, while one independent scale field fixes the conformal factor.
\end{corollary}

\subsection{Connection identification from parallel transport}

\begin{theorem}[Parallel-transport identification of the Recognition connection]
\label{thm:connection-identification}
Let $P\to M$ be a principal $G$-bundle over a connected manifold, and let $A$ and $\widetilde A$ be two smooth connections.  Fix a base point $x_0$ and an identification of the two fibres over $x_0$.  Suppose that, under this identification, the holonomies of $A$ and $\widetilde A$ agree for every piecewise smooth loop based at $x_0$.  Then there exists a smooth gauge transformation $u$ such that
\begin{equation}
\widetilde A=uAu^{-1}-(\dd u)u^{-1}.
\label{eq:gauge-identification}
\end{equation}
Equivalently, complete parallel-transport data identify a connection up to gauge.
\end{theorem}

\begin{proof}
For $x\in M$, choose a path $\gamma_x$ from $x_0$ to $x$ and define $u_x$ by transporting the fixed fibre identification along $\gamma_x$ using $A$ in one direction and $\widetilde A$ in the other.  If another path is chosen, the two definitions differ by the relative holonomy around the resulting based loop, which is the identity by hypothesis.  Thus $u_x$ is path independent.  Smooth dependence of parallel transport on endpoints makes $u$ smooth.  By construction, $u$ intertwines parallel transport along every path.  Differentiating this intertwining along infinitesimal paths gives \eqref{eq:gauge-identification}; see also the connection--transport correspondence in \cite{KobayashiNomizu}.
\end{proof}

\subsection{Recognition Kernel interpretation}

Define the geometric observation
\[
\mathcal O_{\mathrm{geom}}(g,A)
:=\bigl(\mathcal N_g,\mathrm{scale}_g,\operatorname{Hol}_A\bigr),
\]
where $\mathrm{scale}_g$ denotes either the volume density or a pointwise calibration along a declared nowhere-null field.  Let the target be the pair $(g,A)$ modulo diffeomorphisms and gauge transformations.

\begin{corollary}[Target-faithful spacetime identification]
\label{cor:target-faithful-spacetime}
On the class of smooth time-oriented Lorentzian metrics in dimension $n\ge3$ and smooth Recognition connections, equality of $\mathcal O_{\mathrm{geom}}$ implies equality of the target class.  Equivalently, every difference invisible to the observation lies entirely inside a diffeomorphism--gauge orbit.  Consequently, matching characteristic propagation, one scale calibration, and complete transport data identifies the Recognition metric and connection with the observed spacetime geometry up to diffeomorphism and gauge.
\end{corollary}

\begin{proof}
The metric statement is Theorem~\ref{thm:null-cone-identification}; the connection statement is Theorem~\ref{thm:connection-identification}.  Their product says that the observation map is injective on the quotient by the declared target symmetries.
\end{proof}

\begin{remark}[Experimental boundary after identifiability]
The corollary proves the sufficiency and uniqueness of the observation package.  Whether physical measurements actually satisfy those hypotheses is an empirical question.  The remaining experimental task is therefore sharply falsifiable: compare observed wavefront cones, calibrated scale, and parallel transport with the Recognition predictions.  A mismatch is an open residue, not a licence to redefine the target after the experiment.
\end{remark}
